%===============================================================================
% Calibrate Where You Deploy: Conformal Sensor Escalation Under Operating-Point Shift
% IFAC SAFEPROCESS 2027 submission -- ifacconf class (IFAC template, July 2026)
% Build:  python paper/make_tables_ifac.py   (from the repository root)
%         cd paper && pdflatex main_ifac && bibtex main_ifac && pdflatex main_ifac && pdflatex main_ifac
%===============================================================================
\documentclass{ifacconf}
\usepackage{graphicx}
\usepackage{natbib}
\usepackage{amsmath,amssymb}
\usepackage{booktabs}
\usepackage{siunitx}
\usepackage{xcolor}
\newcommand{\covEightyPlr}{50.0}
\newcommand{\covEightyPrf}{32.6}
\newcommand{\gapGroupedK}{+0.045}
\newcommand{\gapLoloK}{-0.148}
\newcommand{\gapLoloWorstK}{-0.453}
\newcommand{\gapGroupedP}{-0.009}
\newcommand{\gapLoloP}{-0.296}
\newcommand{\gapLoloWorstP}{-0.501}
\newcommand{\gapBearingP}{+0.002}
\newcommand{\gateDeltaK}{-0.0008 \pm 0.0051}
\newcommand{\gateDeltaP}{-0.0028 \pm 0.0052}
\newcommand{\gateDeltaKshift}{-0.0037 \pm 0.0145}
\newcommand{\gateDeltaPshift}{-0.0016 \pm 0.0048}
\newcommand{\recalSrcK}{-0.144}
\newcommand{\recalFiveK}{-0.001}
\newcommand{\recalThreeTK}{-0.052}
\newcommand{\recalThreeSTK}{-0.100}
\newcommand{\recalSrcP}{-0.337}
\newcommand{\recalFiveP}{+0.008}
\newcommand{\recalThreeTP}{-0.037}
\newcommand{\recalThreeSTP}{-0.327}
\newcommand{\recalEightSTP}{-0.311}
\newcommand{\energyRatio}{25}
\newcommand{\damageFone}{0.11--0.35}
\newcommand{\escCurG}{33}
\newcommand{\fCurG}{0.939}
\newcommand{\escCurL}{53}
\newcommand{\fCurL}{0.722}
\newcommand{\memsMj}{0.53}
\newcommand{\energyCurGFree}{0.18}
\newcommand{\mcuMemsMj}{2.6}
\newcommand{\mcuRatio}{5}
\newcommand{\covEightyKlr}{40.2}
\newcommand{\covEightyKrf}{70.3}
\newcommand{\computeRatio}{7}
\newcommand{\srcCalP}{7\,729}
\newcommand{\ratioTmcs}{7}
\newcommand{\computeTot}{0.071}
\newcommand{\cycleRatio}{1.9}
\newcommand{\waitiMj}{17.4}
\newcommand{\wFivePthree}{-0.271}
\newcommand{\wFivePlow}{-0.056}
\newcommand{\wFivePtgt}{+0.008}
\newcommand{\wFiveKthree}{-0.030}
\newcommand{\wFiveKtgt}{-0.001}
\newcommand{\escInPrf}{12}
\newcommand{\escTgtPrf}{68}
\newcommand{\escInPlr}{68}
\newcommand{\escTgtPlr}{91}
\newcommand{\sizeInPrf}{1.1}
\newcommand{\sizeTgtPrf}{2.1}
\newcommand{\escSrcMeanK}{75}
\newcommand{\escTgtMeanK}{92}
\newcommand{\escSrcMeanP}{42}
\newcommand{\escTgtMeanP}{72}
\newcommand{\spreadK}{0.16}
\newcommand{\spreadP}{0.13}
\newcommand{\belowK}{37}
\newcommand{\belowP}{28}
\newcommand{\wrongLrK}{0.44}
\newcommand{\wrongLrP}{0.46}
\newcommand{\wrongRfK}{0.06}
\newcommand{\wrongRfP}{0.05}
\newcommand{\escCurLlo}{26}
\newcommand{\escCurLhi}{70}
\newcommand{\ppWorstK}{-0.743}
\newcommand{\ppWorstP}{-0.825}
\newcommand{\mlpDeficit}{0.014}
\newcommand{\fCurGrpP}{0.91}
\newcommand{\fCurBearP}{0.46}
\newcommand{\fCurShiftP}{0.61}
\newcommand{\escVibG}{6}
\newcommand{\fVibG}{0.890}
\newcommand{\escVibL}{6}
\newcommand{\fVibL}{0.719}
\newcommand{\escVibLrG}{79}
\newcommand{\escVibLrL}{69}
\newcommand{\escCurLrG}{88}
\newcommand{\escCurLrL}{49}
\newcommand{\wrongVibRfPten}{0.28}
\newcommand{\wrongCurRfPten}{0.20}
\newcommand{\fullsetK}{15/18}
\newcommand{\memsTotMj}{3.1}
\newcommand{\energyCurGFreeTot}{1.0}
\newcommand{\breakEvenS}{4}

\graphicspath{{figs/}{../figs_en/}}
\newcommand{\TODO}[1]{\textcolor{red}{[#1]}}

\begin{document}
\begin{frontmatter}

\title{Calibrate Where You Deploy: Conformal Sensor Escalation Under Operating-Point Shift\thanksref{footnoteinfo}}

\thanks[footnoteinfo]{\TODO{Funding acknowledgement.}}

\author[USP]{C\^andido A. C. de Lucena Filho}
\author[USP]{Leonardo Fajardo Grupioni}
\author[USP]{Felipe Valencia de Almeida}

\address[USP]{Polytechnic School of the Universidade de S\~ao Paulo, S\~ao Paulo, SP, Brazil (e-mail: candidoalfredo@usp.br, leogrupioni@usp.br, fvalencia@usp.br).}

\begin{abstract}
Condition-monitoring nodes keep a cheap sensor on and wake an expensive one on demand. We replace the tuned threshold governing this escalation with a split-conformal gate: the risk level becomes the specification. On the KAIST and Paderborn datasets the gate matches a tuned threshold in accuracy, but its coverage guarantee, valid in distribution, collapses under operating-point shift: nominal 80\,\% delivers 33--70\,\%. Speed shifts break every modality; torque breaks only motor current. Five recordings from the new operating point restore coverage; pooling them with stale calibration data does not. Recalibration must be triggered by the operating point.
\end{abstract}

\begin{keyword}
Fault detection and diagnosis; condition monitoring; conformal prediction; distribution shift; sensor selection; embedded systems; rolling-element bearings.
\end{keyword}

\end{frontmatter}
%===============================================================================

\section{Introduction}

A battery-powered condition-monitoring node is dominated, energetically, by acquisition: keeping a sensing chain powered for a 160\,ms window costs far more than classifying that window (Table~\ref{tab:energy}). The usual answer is a cascade. A cheap \emph{tier} (a set of channels and the features computed from them) is always on, and the expensive tier is woken only when the cheap tier looks uncertain. The weak point of such designs is the wake-up rule. In practice it is a \emph{threshold gate}, $\max_y \hat p(y\mid x)<\tau$ on the top-class score of the cheap classifier, with $\tau$ swept until the accuracy--energy trade-off looks acceptable. Such scores are not calibrated probabilities, so $\tau$ has no operational meaning, and the escalation rate, which sets battery life, is an outcome of tuning rather than a specification.

Split conformal prediction \citep{vovk2005,lei2018} returns prediction \emph{sets} with finite-sample marginal coverage $\mathbb P(Y\in C(X))\ge 1-\alpha$ under exchangeability and with no assumption on the classifier. Using the set as the gate (escalate iff $|C(x)|\ne 1$) replaces $\tau$ by a stated risk $\alpha$ and makes the escalation rate a measurement. This rule has recently been applied to cascades of \emph{models}, routing images between light and heavy networks \citep{turkoglu2026cads} or queries between tiers of language models \citep{dou2026conformalcascade} by prediction-set size; there the cascaded object is a model and the currency is computation. Here the cascaded object is a \emph{sensor}, and the currency is the energy of acquisition.

Two questions decide whether such a gate can be trusted on a machine. First, the guarantee is conditional on exchangeability between calibration and deployment data, which a change of operating point violates \citep{tibshirani2019covariate,barber2023beyond}. Regime-conditional under-coverage has been reported for remaining-useful-life regression \citep{wang2026rulregime}, and cross-domain conformal sets for bearing diagnosis keep nominal coverage only when calibrated on target data \citep{mei2025calibratedset}; calibration at one operating point and deployment at another has not been measured. Second, if the guarantee fails, what does it cost to restore it? Label-complexity results indicate that only target labels can restore coverage under shift \citep{han2026labelcomplexity}, but no commissioning number exists for condition monitoring.

This paper answers both questions on two public benchmarks with one pipeline. (i) The conformal set-size gate is moved from model cascades to sensor wake-up and priced in datasheet energy per window; with a dedicated transducer, motor current is the expensive tier, so a sensor cascade that saves acquisition energy exists only when current is free, and that cascade is the one most exposed to shift. (ii) At matched escalation fraction the conformal gate and a tuned threshold gate are indistinguishable in accuracy; the value of the gate is the specification, not the predictor. (iii) With source-only calibration, nominal 80\,\% coverage falls to \covEightyKlr--\covEightyKrf\,\% (KAIST) and \covEightyPrf--\covEightyPlr\,\% (Paderborn) depending on the classifier; the gate then fails in one of two modes that the node cannot tell apart, and on Paderborn the collapse is modality-specific: speed breaks every tier, torque only the current tier, radial force the vibration tiers moderately. (iv) Five labelled recordings from the new operating point restore mean coverage to within 0.01 of nominal (three to within about 0.05); pooling them with the stale source scores is worse than replacing them, and weighted conformal prediction recovers the same effect only when the target share of the normalised weight is large.

\section{Conformal Sensor Escalation}

\subsection{Split conformal prediction}
Let $\hat p(\cdot\mid x)$ be any classifier over $K$ classes and $\{(x_i,y_i)\}_{i=1}^n$ a calibration set disjoint from training. With the least-ambiguous-set score (LAC) $s(x,y)=1-\hat p(y\mid x)$ \citep{sadinle2019} and $\hat q$ the $\lceil (n+1)(1-\alpha)\rceil$-th smallest calibration score, the prediction set is
\begin{equation}
C(x)=\{\,y: s(x,y)\le \hat q\,\}.
\label{eq:set}
\end{equation}
If calibration and test points are exchangeable, $\mathbb P(Y\in C(X))\ge 1-\alpha$ for any finite $n$ \citep{vovk2005,lei2018}. The adaptive-prediction-set score of \citet{romano2020} (APS) is also evaluated; it buys conditional coverage with larger sets. Under weighted conformal prediction \citep{tibshirani2019covariate,barber2023beyond} calibration score $i$ carries a weight $w_i\ge 0$, the test point weight 1, and $\hat q$ is the $(1-\alpha)$-quantile of the weighted empirical distribution normalised by $\sum_i w_i+1$; with $w_i\equiv 1$ this is (\ref{eq:set}). For fixed weights the applicable guarantee is that of \citet{barber2023beyond}: coverage at least $1-\alpha-\sum_i\tilde w_i\,d^{(i)}_{\mathrm{TV}}$, where $\tilde w_i$ are the normalised weights and $d^{(i)}_{\mathrm{TV}}$ the total-variation distance between the data sequence and the one with calibration point $i$ and the test point swapped, so down-weighting the stale scores bounds the loss. One more caveat applies throughout: the windows of one recording are dependent, so the calibration scores are exchangeable at recording level, not window level. The finite-sample statement then holds approximately, with an effective calibration size equal to the number of recordings \citep{dunn2023hierarchical,lee2023hierarchical}; all protocols below therefore reason in recordings, and small $\alpha$ cannot be certified from three of them.

\subsection{The gate}\label{sec:thegate}
The cheap tier runs on every window; the expensive tier is woken, and its classifier decides, exactly when the cheap tier's set is not a singleton:
\begin{equation}
\text{escalate}(x)\iff |C_{\text{cheap}}(x)|\ne 1.
\label{eq:gate}
\end{equation}
Anything but a singleton (several classes, or none) is an admission that, at risk $\alpha$, the cheap sensor cannot resolve the window. The gate does not promise that the singleton is correct, (\ref{eq:set}) being marginal; it promises that the fraction of all windows on which the node emits a wrong singleton is at most $\alpha$ (the error rate \emph{among} singletons is not bounded), and it turns the escalation fraction, hence the duty cycle of the expensive sensor, into an observable. With the LAC score, $|C(x)|=1$ iff the top-class probability is at least $1-\hat q$ and the second is below it. When $1-\hat q>0.5$ this is a threshold on the top-class score, which is why the threshold gate is the right baseline; but in the experiments $\hat q$ is large ($\hat q\ge 0.80$ for the KAIST current tier at every $\alpha\le 0.20$ and $\ge 0.65$ at $\alpha\le 0.40$; on Paderborn $1-\hat q>0.5$ only for the random forest at $\alpha\ge 0.20$), so the two gates are different selection rules and their comparison in Sec.~\ref{sec:gate} is empirical. In model-based terms the operating point is a measured scheduling variable the classifier does not see; conditioning the calibration on an operating-point tag (Mondrian conformal prediction, \citealp{vovk2005}) is the data-driven counterpart of a scheduled residual threshold, and Sec.~\ref{sec:recal} bounds the data it costs per operating point.

\subsection{Cost: operations and energy}
Two currencies are kept separate. Operations per window follow from the tier: $9N$ per channel for the nine time-domain statistics, plus $\approx 2.5N\log_2 N+2N$ with the FFT block ($N=4096$); they are identical across datasets. Energy is modelled from datasheets (Table~\ref{tab:energy}): the chain of a tier is powered for the window, so its energy is power $\times$ 160\,ms, and arithmetic is priced in picojoules per cycle times published cycle counts, marginally and on a total-power basis that includes the idle floor a duty-cycled node pays. An active current transducer draws \ratioTmcs--\energyRatio$\times$ the power of a three-axis digital MEMS accelerometer per phase (Hall sensor to isolated transducer), so with dedicated active sensing the ``cheap current'' tier is the expensive one in joules. The computation of the accelerometer time+FFT tier (\computeTot\,mJ on a total basis) is \computeRatio$\times$ below the energy of keeping the MEMS sensor on, so a cascade between two tiers of the \emph{same} sensor saves operations (4.6$\times$; \cycleRatio$\times$ in cycles, since the FFT runs in SIMD assembly), not acquisition energy. The only cascade that saves acquisition energy, current always on and vibration on demand, requires current to be free, i.e.\ the node is the drive controller that already samples phase current (at a lower rate and resolution than the laboratory signal, so the current-tier classifier must be trained on inverter-sampled data); its energy per window is then the MEMS energy times the escalation fraction. The MCU state \emph{during} acquisition is not free either: the IIS3DWB FIFO holds an eighth of a window, so the ESP32-S3 sleeps between eight FIFO drains over SPI and spends \mcuMemsMj\,mJ, \mcuRatio$\times$ the sensor itself; the current ADC has no FIFO, so on a dedicated current tier the MCU idles for the whole window (\waitiMj\,mJ). At one window per minute the MEMS tier with its MCU service averages \memsAvgUwTot\,\si{\micro\watt} against a \deepSleepUw\,\si{\micro\watt} deep-sleep floor, so the joules of the cascade matter for near-continuous monitoring; sensor turn-on time is excluded.

\begin{table}[tb]\centering\footnotesize
\caption{Datasheet power of the sensing chains and of the MCU, and energy per 160\,ms window. The MEMS FIFO holds an eighth of a window, so the MCU drains it eight times over SPI (about 20\,ms at the idle floor); the current ADC has no FIFO. Last row: ESP32-S3 computation for one window of the accelerometer time+FFT tier (71\,264 cycles, FFT count extrapolated from the ESP-DSP benchmarks), marginal at 543\,pJ per cycle and on a total-power basis.}
\label{tab:energy}
\setlength{\tabcolsep}{3.5pt}\resizebox{\columnwidth}{!}{\begin{tabular}{lrr}\hline
component & mW & mJ / window \\ \hline
MEMS accel.\ IIS3DWB, 3 axes, ADC on chip & 3.30 & 0.53 \\
Current transducer LEM CKSR 15-NP + ADC, per phase & 81.66 & 13.07 \\
Hall current sensor TMCS1100 + ADC, per phase & 22.52 & 3.60 \\
IEPE accel.\ on TI TIDA-010249 front end & 100.00 & 16.00 \\
ESP32-S3, MEMS window: light sleep + 8 FIFO drains & -- & 2.25 \\
ESP32-S3, current window (no FIFO): idle, 240\,MHz & 108.57 & 17.37 \\
ESP32-S3 computation, accel.\ time+FFT tier (marginal / total) & -- & 0.039 / 0.071 \\
\hline\end{tabular}}\end{table}

\section{Experimental Setup}

\textbf{Data.} KAIST \citep{jung2023kaist}: a motor-driven shaft with four accelerometers, three current transformers (Hioki CT6700) and two housing temperatures, vibration and current both at 25.6\,kHz; five conditions at three torque loads (0, 2, 4\,N\,m): normal (one recording per load), inner- and outer-race faults, misalignment and unbalance at three to five severities each, pooled per class; 45 recordings; only the U-phase current is present in every file and is the one used. Paderborn \citep{lessmeier2016paderborn}: 32 bearings (6 healthy, 12 with machined damage, 14 with damage from accelerated-lifetime tests), one vibration channel and two current phases at 64\,kHz, four operating points, three of which differ from the baseline (1500\,rpm, 0.7\,N\,m, 1000\,N) in exactly one factor: speed 900\,rpm, torque 0.1\,N\,m, radial force 400\,N; four classes (healthy, outer race, inner race, combined; the combined class has three bearings, all with real damage). Signals are resampled to 25.6\,kHz with an exact 2/5 polyphase filter, which discards the resonance band above 12.8\,kHz used by envelope analysis; five of the twenty measurements per bearing and operating point are used (639 of 640 recordings, 31\,266 windows).

\textbf{Features and tiers.} Windows of 4096 samples with hop 2048, at most 50 per recording; per channel nine time-domain statistics (mean, standard deviation, RMS, peak-to-peak, kurtosis, skewness, crest, shape and impulse factors) and, when the FFT block is included, six spectral descriptors (dominant frequency, centroid, spectral entropy, energy in 0--2, 2--6 and 6--12.8\,kHz); identical code for both datasets.

\textbf{Classifiers.} $\ell_2$-regularised logistic regression (LR, $C=1$ on standardised features) and a random forest (RF, 300 trees, minimum leaf 2); convergence is checked for every fit. A 64--32 MLP, never the best classifier, is omitted from the tables.

\textbf{Validation.} All splits are at recording level. \emph{Recording-grouped} stratified 3-fold cross-validation (three seeds) estimates in-distribution performance (``in'' in the tables); \emph{operating-point shift} trains on all but one operating point and tests on it (leave-one-operating-point-out). On Paderborn the same bearing is recorded at every operating point, so recording-grouped folds let a bearing appear on both sides; a \emph{bearing-grouped} protocol (folds and calibration sets carved by bearing) bounds that optimism, with the caveat that the three combined-damage bearings fall one per fold. The shift protocol shares bearings by construction, which attributes the degradation it exposes to the operating point alone. Intervals are 95\,\% percentile intervals of a cluster bootstrap over recordings (600 replicates).

\textbf{Conformal protocol.} A stratified third of the training recordings is held out for calibration in every fold; the calibration scores of the three folds are pooled and one $\hat q$ is applied to the pooled test predictions, so that under shift the pooled calibration set holds scores from every operating point through the other folds' classifiers; the recalibration protocol below uses one quantile per held-out point. Coverage is the fraction of test windows whose label lies in $C(x)$, and the \emph{gap} is coverage minus $1-\alpha$, for $\alpha\in\{0.01,0.02,0.05,0.10,0.15,0.20,0.30,0.40\}$. The conformal gate is compared with the threshold gate at matched escalation fraction by interpolating the swept threshold curve, in distribution. For recalibration each held-out operating point is the target: the classifier is trained on the others and the calibration set is the source recordings only, the source recordings pooled with $k$ labelled target recordings, or the $k$ target recordings alone, with $k\in\{1,2,3,5,8\}$ drawn at random over five seeds and $\alpha\in\{0.05,0.10,0.20,0.30\}$; test recordings are disjoint from calibration recordings. Weighted conformal prediction gives the source scores a fixed weight $\rho\in\{1,0.3,0.1,0.03,0\}$ and the target scores weight 1.

\section{Results}\label{sec:results}

\begin{table*}[tb]\centering\footnotesize
\caption{Cost--accuracy frontier: macro-F1 of the best classifier per tier (mean over three seeds for the grouped protocols; one run per held-out point under shift) with 95\,\% percentile intervals of a cluster bootstrap over recordings (bearings for the bearing-grouped column) of the seed-0 predictions; the interval can sit off the point where a class has few recordings (three per class on KAIST). Operations per 160\,ms window.}
\label{tab:frontier}
\setlength{\tabcolsep}{3.5pt}\resizebox{0.9\textwidth}{!}{\begin{tabular}{llrlll}\hline
 & tier & ops / window & recording-grouped [95\,\% CI] & bearing-grouped [95\,\% CI] & operating-point shift [95\,\% CI] \\ \hline
KAIST & current, time+FFT & 167\,936 & 0.417 [0.33, 0.45] & -- & 0.213 [0.13, 0.29] \\
 & 1 accel., time & 36\,864 & 0.690 [0.66, 0.76] & -- & 0.687 [0.63, 0.74] \\
 & 1 accel., time+FFT & 167\,936 & 0.706 [0.68, 0.78] & -- & 0.714 [0.61, 0.80] \\
 & 4 accel., time+FFT & 671\,744 & 0.729 [0.70, 0.91] & -- & 0.746 [0.68, 0.81] \\
 & 4 accel. + current & 839\,680 & 0.729 [0.70, 0.92] & -- & 0.675 [0.60, 0.74] \\
Paderborn & 2 current phases, time+FFT & 335\,872 & 0.908 [0.89, 0.92] & 0.458 [0.34, 0.61] & 0.612 [0.58, 0.64] \\
 & 1 accel., time & 36\,864 & 0.888 [0.88, 0.91] & 0.548 [0.41, 0.68] & 0.714 [0.68, 0.75] \\
 & 1 accel., time+FFT & 167\,936 & 0.959 [0.95, 0.97] & 0.553 [0.40, 0.67] & 0.762 [0.73, 0.79] \\
 & 1 accel. + 2 phases & 503\,808 & 0.991 [0.98, 0.99] & 0.504 [0.37, 0.63] & 0.768 [0.74, 0.80] \\
\hline\end{tabular}}\end{table*}

\subsection{How much sensing is needed}
On KAIST (Table~\ref{tab:frontier}), nine time-domain features from a single accelerometer (36\,864 operations) reach 0.69 macro-F1 in distribution and under shift, against 0.73/0.75 for four accelerometers with FFT (671\,744 operations), and current alone collapses under shift (0.42 to 0.21). On Paderborn, two current phases reach 0.91 under recording-grouped folds, nearly on par with vibration, but the bearing-grouped column shows that recording-grouped folds measure recognition of known bearings rather than of damage classes on unseen bearings: every tier falls to 0.46--0.55 \citep[cf.][]{vieira2025realistic}. Under operating-point shift vibration degrades (0.959 to 0.762) and current collapses (0.908 to 0.612), and fusion does not help where it matters (0.991 to 0.768). Two cascades are evaluated, each with the full vibration tier of its dataset as the expensive tier (four accelerometers with FFT on KAIST, one on Paderborn): current-first (\emph{current} $\to$ \emph{vibration, time+FFT}), the only one that can save acquisition energy, and same-sensor (\emph{one accelerometer, time-domain} $\to$ \emph{vibration, time+FFT}), which saves computation.

\subsection{The gate does not buy accuracy}\label{sec:gate}
The conformal gate is compared with the threshold gate at matched escalation fraction, i.e.\ at equal cost. In distribution, pooled over cascades, scores, classifiers and seeds, the macro-F1 difference (mean $\pm$ s.d.\ over matched points) is $\gateDeltaK$ on KAIST and $\gateDeltaP$ on Paderborn, and between $-0.003$ and $+0.001$ for every cascade and score. The two gates trace the same accuracy--escalation curve although, as noted in Sec.~\ref{sec:thegate}, they are different selection rules at the $\hat q$ observed. This negative result is a counterpoint to the motivation of \citet{dou2026conformalcascade}: the conformal gate is not a better predictor than a tuned threshold; what it changes is that $\alpha$ fixes the operating point without tuning.

\begin{table}[tb]\centering\footnotesize
\caption{Coverage gap (empirical $-$ nominal) of LAC sets for the current tier (current, time+FFT on KAIST; two phases on Paderborn): mean and worst cell over the $\alpha$ grid, classifiers and seeds (48 cells per row); coverage at nominal 80\,\% by LR / RF, mean $\pm$ s.d.\ over seeds, and the same averaged over classifiers for the healthy and the fault classes. in = recording-grouped, bearing = bearing-grouped, shift = operating-point shift.}
\label{tab:gap}
\setlength{\tabcolsep}{2pt}\resizebox{\columnwidth}{!}{\begin{tabular}{@{}llcccc@{}}\hline
dataset & protocol & mean & worst & LR / RF at 80\,\% & healthy / faulty \\ \hline
KAIST & in & $+0.045$ & $-0.019$ & 82.9$\pm$2.1 / 93.5$\pm$11.2& 38 / 92 \\
KAIST & shift & $-0.148$ & $-0.453$ & 40.2$\pm$5.4 / 70.3$\pm$6.4& 33 / 56 \\
Paderborn & in & $-0.009$ & $-0.039$ & 79.0$\pm$1.3 / 78.2$\pm$0.9& 69 / 82 \\
Paderborn & bearing & $+0.002$ & $-0.053$ & 81.2$\pm$4.7 / 80.8$\pm$1.0 & 62 / 82 \\
Paderborn & shift & $-0.296$ & $-0.501$ & 50.0$\pm$5.9 / 32.6$\pm$2.6& 57 / 37 \\
\hline\end{tabular}}\end{table}

\subsection{Coverage collapses under operating-point shift}
Table~\ref{tab:gap} gives the central result for the current tier, the cheap tier of the energy-saving cascade; Table~\ref{tab:factor} extends it to the vibration tiers. Under recording-grouped validation the marginal guarantee holds (mean gap $\gapGroupedK$ on KAIST, $\gapGroupedP$ on Paderborn), and under bearing-grouped validation as well ($\gapBearingP$ mean, worst $-0.053$): marginal coverage does not depend on the accuracy of the classifier, the property the gate relies on. Under operating-point shift it fails: mean gap $\gapLoloK$ (worst $\gapLoloWorstK$) on KAIST and $\gapLoloP$ (worst $\gapLoloWorstP$) on Paderborn; at nominal 80\,\% the gate delivers \covEightyKlr--\covEightyKrf\,\% on KAIST and \covEightyPrf--\covEightyPlr\,\% on Paderborn (LR to RF on KAIST, RF to LR on Paderborn). The APS score stays near nominal only by inflating sets towards the full label set, and the failure is uneven across operating points (repository notebooks). The healthy/faulty column is the diagnostic reading of the same numbers: marginal coverage says nothing about the healthy class, which on the KAIST current tier is covered at 38\,\% in distribution (9\,\% LR, 67\,\% RF; one healthy recording per load) while the fault classes are covered at 92\,\%; a false-alarm bound needs class-conditional (Mondrian) calibration, with the data cost that Sec.~\ref{sec:recal} measures.

How the gate fails depends on how the cheap classifier behaves off-distribution (Table~\ref{tab:factor}, right), and only one of the two modes is silent. Logistic regression keeps emitting singletons under shift (0.62--0.96 of the sets) while coverage falls to 0.46--0.68: the sets are confident \emph{and wrong}, so the gate under-escalates. The random forest instead returns empty sets (0.47--0.75): a quantile calibrated on confident in-distribution probabilities is small, under shift no class clears it, and by (\ref{eq:gate}) the gate over-escalates, which is wasteful but safe. Which mode occurs depends on the classifier's off-distribution confidence and on the shift (trained on three operating points, Table~\ref{tab:shift}, the random forest of the same-sensor cascade keeps its escalation fraction and loses accuracy silently); the node cannot observe it. Its direction is modality-specific: under a speed change the LR vibration gates still cover the healthy class fully and the fault classes at 0.03--0.55 (missed detections), whereas the LR current gate under speed or torque change covers the healthy class at 0.00 (false alarms).


\begin{table*}[tb]\centering\footnotesize
\caption{Paderborn, LAC at nominal 80\,\%, classifiers fitted and calibrated at the baseline operating point: coverage (macro-F1 in parentheses) on held-out baseline recordings (in) and on each operating point that changes one factor; coverage, singleton and empty-set fractions with the three shifted points pooled. A wrong singleton is silent under-escalation; an empty set escalates by (\ref{eq:gate}).}
\label{tab:factor}
\setlength{\tabcolsep}{3.5pt}\resizebox{0.9\textwidth}{!}{\begin{tabular}{llcccc ccc}\hline
 & & \multicolumn{4}{c}{coverage (macro-F1) by factor changed} & \multicolumn{3}{c}{shift, pooled} \\
tier & classifier & in & speed & torque & radial force & coverage & singletons & empty \\ \hline
2 current phases, time+FFT & LR & 0.77 (0.74) & 0.38 (0.14) & 0.35 (0.13) & 0.66 (0.57) & 0.46 & 0.96 & 0.00 \\
 & RF & 0.78 (0.92) & 0.01 (0.26) & 0.15 (0.43) & 0.46 (0.66) & 0.21 & 0.25 & 0.75 \\
1 accel., time & LR & 0.82 (0.66) & 0.50 (0.31) & 0.82 (0.64) & 0.73 (0.56) & 0.68 & 0.62 & 0.00 \\
 & RF & 0.79 (0.93) & 0.08 (0.22) & 0.71 (0.86) & 0.40 (0.59) & 0.40 & 0.53 & 0.47 \\
1 accel., time+FFT & LR & 0.83 (0.80) & 0.36 (0.27) & 0.80 (0.76) & 0.64 (0.63) & 0.60 & 0.91 & 0.01 \\
 & RF & 0.86 (0.98) & 0.07 (0.26) & 0.67 (0.90) & 0.27 (0.62) & 0.34 & 0.39 & 0.61 \\
\hline\end{tabular}}\end{table*}

\subsection{Which factor breaks exchangeability, per modality}
Because three Paderborn operating points each change one factor, classifiers fitted and calibrated at the baseline can be evaluated on each neighbour (Table~\ref{tab:factor}, left; coverage below nominal 0.80 is the reading). For the vibration tiers a change of speed destroys coverage (0.07--0.50), a change of torque leaves it near nominal for LR (0.80--0.82) and 0.09--0.13 below nominal for RF, and a change of radial force degrades it moderately (0.64--0.73 for LR, 0.27--0.40 for RF); for the current tier both speed and torque destroy it (0.01--0.38). Macro-F1 tells the same story (vibration under torque change 0.64--0.90, under speed change 0.22--0.31). This agrees with accuracy-based domain-shift studies on this benchmark, which find that speed drives the degradation while load and radial force have small effects \citep{panic2027baselines}: those studies use vibration, and the torque sensitivity belongs to the current tier, as the KAIST shift column (vibration 0.69 to 0.69, current 0.42 to 0.21) independently confirms. The physics is consistent. Speed scales impact rate and vibration amplitude, so the amplitude statistics move and the scale-free ones (kurtosis, crest factor) move through the number of impacts per window, and on an inverter-fed motor it moves the current fundamental and every sideband. Torque from 0.7 to 0.1\,N\,m is a roughly seven-fold change in q-axis current: every amplitude statistic of the current channel rescales while the fault sidebands stay small relative to the fundamental, the classical load dependence of motor-current signature analysis, which the gate inherits; vibration is mainly unaffected. Radial force changes the load zone, hence the amplitude of outer-race impulses. The design consequence: a current-first cascade must be recalibrated whenever speed \emph{or} torque changes; a vibration tier tolerates torque changes and needs recalibration for speed and, to a lesser degree, radial-force changes.

\begin{table}[tb]\centering\footnotesize
\caption{Paderborn: the LAC gate at a fixed $\alpha$ in distribution (recording-grouped) and under operating-point shift, mean over seeds. esc = fraction of windows escalated; F1 = macro-F1 of the cascade; exp.\ = macro-F1 of the expensive tier always on, fitted without the calibration recordings (hence below Table~\ref{tab:frontier}).}
\label{tab:shift}
\setlength{\tabcolsep}{3.5pt}\resizebox{\columnwidth}{!}{\begin{tabular}{llc ccc ccc}\hline
 & & & \multicolumn{3}{c}{in distribution} & \multicolumn{3}{c}{operating-point shift} \\
cascade & classifier & $\alpha$ & esc & F1 & exp. & esc & F1 & exp. \\ \hline
1 accel.\ time $\to$ 1 accel.+FFT & LR & 0.10 & 79\,\% & 0.688 & 0.687 & 69\,\% & 0.614 & 0.612 \\
 & RF & 0.10 & 6\,\% & 0.890 & 0.941 & 6\,\% & 0.719 & 0.753 \\
2 phases $\to$ 1 accel.+FFT & LR & 0.05 & 88\,\% & 0.753 & 0.687 & 49\,\% & 0.532 & 0.612 \\
 & RF & 0.05 & 33\,\% & 0.939 & 0.941 & 53\,\% & 0.722 & 0.753 \\
\hline\end{tabular}}\end{table}

\begin{table}[tb]\centering\footnotesize
\caption{Recalibration at the target operating point: coverage gap (LAC) by number $k$ of target recordings in the calibration set and by composition (source recordings pooled with the $k$ target recordings; the $k$ target recordings only), mean over classifiers, $\alpha\in\{0.05,0.10,0.20,0.30\}$ and held-out points, $\pm$ s.d.\ over five seeds (random draw of the $k$ recordings, not stratified by class). Source-only calibration gives $-0.144\pm0.038$ (KAIST) and $-0.337\pm0.021$ (Paderborn).}
\label{tab:recal}
\setlength{\tabcolsep}{3.5pt}\resizebox{\columnwidth}{!}{\begin{tabular}{c cc c cc}\hline
 & \multicolumn{2}{c}{KAIST} & & \multicolumn{2}{c}{Paderborn} \\
$k$ & pooled & target & & pooled & target \\ \hline
1 & $-0.153\pm0.034$ & $-0.240\pm0.058$ & & $-0.334\pm0.019$ & $-0.270\pm0.119$ \\
2 & $-0.127\pm0.027$ & $-0.081\pm0.045$ & & $-0.332\pm0.019$ & $-0.125\pm0.076$ \\
3 & $-0.100\pm0.024$ & $-0.052\pm0.051$ & & $-0.327\pm0.018$ & $-0.037\pm0.051$ \\
5 & $-0.074\pm0.048$ & $-0.001\pm0.071$ & & $-0.320\pm0.019$ & $+0.008\pm0.043$ \\
8 & $-0.038\pm0.027$ & $+0.010\pm0.037$ & & $-0.311\pm0.020$ & $+0.014\pm0.039$ \\
\hline\end{tabular}}\end{table}

\subsection{What it takes to fix it}\label{sec:recal}
Table~\ref{tab:recal} quantifies recalibration against the number $k$ of recordings available from the target operating point. Five recordings restore mean coverage to within 0.01 of nominal (gap from $\recalSrcK$ to $\recalFiveK$ on KAIST and from $\recalSrcP$ to $\recalFiveP$ on Paderborn) and three to within about 0.05 ($\recalThreeTK$, $\recalThreeTP$; seed s.d.\ 0.04--0.07, unstratified draw). The number is in \emph{labelled} recordings: on a rig or in a fleet where damaged units can be cycled through the new operating point this is cheap; on a production machine only healthy recordings exist at the new point, which can recalibrate the healthy-class quantile alone (the false-alarm side of Table~\ref{tab:gap}; untested here). Pooling the target recordings with the stale source set is worse than replacing it: at $k=3$ the target-only gap on KAIST is $\recalThreeTK$ against $\recalThreeSTK$ pooled; on Paderborn the source set holds about \srcCalP{} windows against about 50 per target recording, so pooling barely moves the quantile ($\recalThreeSTP$ at $k=3$, $\recalEightSTP$ at $k=8$) while target-only reaches $\recalThreeTP$. Stale scores dominate the quantile by volume; the guarantee needs exchangeability, not volume.

Weighted conformal prediction, with the source scores weighted $\rho$ and the target scores 1 (full grid in the repository), is the principled alternative to discarding the source scores. Its effect follows the target share of the normalised weight, $n_{\text{tgt}}/(n_{\text{tgt}}+\rho\,n_{\text{src}})$: on KAIST ($n_{\text{src}}=500$) $\rho=0.3$ at $k=5$ gives $\wFiveKthree$ against $\wFiveKtgt$ target-only; on Paderborn ($n_{\text{src}}\approx\srcCalP$) $\rho=0.3$ leaves the source in control ($\wFivePthree$ at $k=5$) and only $\rho=0.03$ approaches target-only ($\wFivePlow$ against $\wFivePtgt$). Once three target recordings exist, calibrate on them.

\subsection{The gate under shift}\label{sec:gatefail}
Table~\ref{tab:shift} shows what the expired guarantee means for the node. In the same-sensor cascade the RF gate at $\alpha=0.10$ escalates 6\,\% of windows in distribution and under shift alike, but macro-F1 falls from 0.890 to 0.719; in the current-first cascade with free current, $\alpha=0.05$ with RF escalates \escCurG\,\% of windows in distribution (\energyCurGFree\,mJ of MEMS sensor energy per window instead of \memsMj, MCU service excluded) for \fCurG, and \escCurL\,\% under shift for \fCurL. An operating-point change moves the node along the energy--accuracy curve with no change in $\alpha$ and no signal the node can read.

\subsection{Limits of recalibration}
Recalibration restores the guarantee only when the classifier remains informative at the deployment point. Training on bearings with machined damage and testing on bearings with real damage (and vice versa) gives macro-F1 of \damageFone{} for every tier and classifier (repository notebook, experiment X3); the machined set has no combined-damage class, so this is a three-class model tested on four classes and macro-F1 is at most 0.75 by construction, yet the result is at chance. Target-only recalibration then moves coverage towards nominal only by inflating the sets towards the full label set, so the gate would escalate every window. Operating-point shift is a calibration problem that a few labelled recordings fix; damage-origin shift is a representation problem that they do not.

\section{Discussion}

A conformal escalation gate is a better \emph{interface} than a threshold gate, not a better predictor, and the guarantee that makes it worth having is exactly what expires first when the machine leaves the operating point it was calibrated at. First, conformal gates should ship with a recalibration procedure, and its size is now known: three to five labelled recordings at the new operating point. Second, the current tier is the fragile one; a variable-speed drive cannot run either cascade without per-speed calibration or speed normalisation, both untested here, and the wake-on-threshold vibration cascade of commercial nodes is expected to share the speed sensitivity. Third, neither failure mode is visible to the node, so the recalibration trigger must be external: an operating-point tag from the drive (what a Mondrian calibration would condition on) or a drift monitor.

\textbf{Limitations.} Energy figures are datasheet-derived, the FFT cycle count is extrapolated (a cycle-count measurement on the ESP32-S3 would replace it), and both tiers share one window (no re-acquisition on escalation). Both benchmarks use hand-crafted features with shallow classifiers; the coverage statements are classifier-agnostic, but the transfer failures may differ for learned or speed-normalised representations. Only labelled target data and fixed source weights were tested; label-free recalibration from unlabelled target windows \citep{yilmaz2022testtime}, likelihood-ratio weighting and online adaptive conformal inference \citep{gibbs2021adaptive} remain to be compared.

\textbf{Code and data.} Both datasets are public; feature extraction, experiments, table generator and executed notebooks (runnable on Google Colab) are at \TODO{repository URL, commit hash}.

\section{Conclusion}
Sensor escalation was formulated as a conformal set-size test: a stated risk level replaces a tuned threshold and the escalation rate becomes a measurement. On two benchmarks the gate matches a tuned threshold in accuracy, and its coverage guarantee, valid in distribution even for weak classifiers, fails under operating-point shift, delivering 33--70\,\% where 80\,\% was requested. Five recordings from the new operating point restore the guarantee; pooling them with the stale calibration set is worse than replacing it. Conformal calibration is a commissioning step tied to the operating point, not a one-off performed in the laboratory.

\section*{Declaration of generative AI and AI-assisted technologies in the writing process}
During the preparation of this work the authors used Claude (Anthropic) in order to improve the readability and language of the manuscript. After using this tool, the authors reviewed and edited the content as needed and take full responsibility for the content of the publication.

\bibliography{refs_ifac}

\end{document}
